\section{Silent Goddess}

\subsection*{Deskripsi Soal}

Di lembar terakhir catatan risetnya, Marisa menemukan salinan artikel
investigasi mengenai sosok dewi misterius bersayap tunggal yang terpergok
menyusup ke fasilitas di kedalaman Hutan Bambu. Potongan artikel berita
tersebut kini tidak dapat dibaca karena telah dienkripsi menggunakan
\textbf{Hill \textit{Cipher}}.

Sang jurnalis yang menyusunnya enggan memberikan matriks kunci secara langsung.
Ia hanya membocorkan satu petunjuk bahwa naskah investigasi tersebut diawali
dengan kata \textbf{``EXPLOSIVE SCOOP''}.

Bantulah Marisa merekonstruksi matriks kunci menggunakan petunjuk awal
\textit{plaintext} tersebut melalui \textbf{\textit{Known Plaintext Attack}}
(KPA). Hitung invers modulo dari matriks kunci yang ditemukan, lalu gunakan
hasilnya untuk mendekripsi sisa \textit{ciphertext} hingga pesan lengkap
artikel terungkap.

\medskip
\noindent\textbf{Berkas lampiran:}
\begin{itemize}[nosep]
  \item \textit{Attachment}: \url{https://drive.google.com/file/d/1aS17ZuSGAWfpcus4F6HEP3wPVXYIHE7n/view?usp=sharing}
  \item \texttt{md5sum}: \texttt{84e86bf1dafb89ec420c114abd60498c}
  \item \texttt{sha256sum}: \texttt{6113d33c379b4727e120d3146cfc39ab3004842c0252848dc56c621ba7a5d25f}
\end{itemize}

\subsection*{Berkas \textit{Ciphertext}}

\ciphertextfile{../../src/25/h.txt}

\subsection*{Langkah Penyelesaian}

\textit{Hill cipher} mudah dipecahkan dengan \textit{known-plaintext attack}.
Berdasarkan petunjuk yang diberikan, 14 karakter pertama dari \textit{plaintext} adalah \textbf{``EXPLOSIVESCOOP''}. 
Misalkan untuk \textit{Hill cipher} dengan $m = 3$ diketahui:
\begin{itemize}[nosep]
    \item $P = (4, 23, 15) \to C = (21, 13, 20)$ (dari blok ``EXP'' $\to$ ``VNU'')
    \item $P = (11, 14, 18) \to C = (4, 23, 6)$ (dari blok ``LOS'' $\to$ ``EXG'')
    \item $P = (8, 21, 4) \to C = (12, 12, 23)$ (dari blok ``IVE'' $\to$ ``MMX'')
    \item Jadi, $K(4, 23, 15) = (21, 13, 20)$, $K(11, 14, 18) = (4, 23, 6)$, dan $K(8, 21, 4) = (12, 12, 23)$
\end{itemize}

\begin{equation*}
\begin{pmatrix} 21 & 4 & 12 \\ 13 & 23 & 12 \\ 20 & 6 & 23 \end{pmatrix}
= K
\begin{pmatrix} 4 & 11 & 8 \\ 23 & 14 & 21 \\ 15 & 18 & 4 \end{pmatrix} \pmod{26}
\quad \to K = C P^{-1} \pmod{26}
\end{equation*}

\begin{itemize}[nosep]
    \item Matriks balikan dari $P$ adalah $P^{-1} = \begin{pmatrix} 4 & 11 & 8 \\ 23 & 14 & 21 \\ 15 & 18 & 4 \end{pmatrix}^{-1} \pmod{26} = \begin{pmatrix} 8 & 24 & 1 \\ 1 & 0 & 24 \\ 24 & 1 & 25 \end{pmatrix}$
    \item Sehingga
\end{itemize}

\begin{equation*}
K = C P^{-1} \pmod{26} = 
\begin{pmatrix} 21 & 4 & 12 \\ 13 & 23 & 12 \\ 20 & 6 & 23 \end{pmatrix}
\begin{pmatrix} 8 & 24 & 1 \\ 1 & 0 & 24 \\ 24 & 1 & 25 \end{pmatrix} \pmod{26}
= \begin{pmatrix} 18 & 22 & 1 \\ 25 & 12 & 7 \\ 16 & 9 & 11 \end{pmatrix}
\end{equation*}

Selanjutnya, sisa dari pesan \textit{ciphertext} didekripsi secara langsung menggunakan matriks balikan (invers) dari $K$.

\subsection*{Langkah-langkah Dekripsi}

Berikut adalah rincian langkah dekripsi untuk tiap blok karakter (Total 240 blok) menjadi vektor numerik, dikali dengan invers matriks kunci modulo 26, lalu dikembalikan ke bentuk teks:

\begin{longtable}{|c|c|c|c|c|}
\hline
\textbf{Blok} & \textbf{Ciphertext} & \textbf{Vektor} & \textbf{Plaintext} & \textbf{Vektor} \\
\hline
\endfirsthead

\multicolumn{5}{c}%
{{\bfseries \tablename\ \thetable{} -- lanjutan dari halaman sebelumnya}} \\
\hline
\textbf{Blok} & \textbf{Ciphertext} & \textbf{Vektor} & \textbf{Plaintext} & \textbf{Vektor} \\
\hline
\endhead

\hline
\multicolumn{5}{|r|}{{lanjut di halaman berikutnya,}} \\
\hline
\endfoot

\hline
\endlastfoot

  1 & VNU & $(21, 13, 20)^T$ & EXP & $(4, 23, 15)^T$ \\
  \hline
  2 & EXG & $(4, 23, 6)^T$ & LOS & $(11, 14, 18)^T$ \\
  \hline
  3 & MMX & $(12, 12, 23)^T$ & IVE & $(8, 21, 4)^T$ \\
  \hline
  4 & SAS & $(18, 0, 18)^T$ & SCO & $(18, 2, 14)^T$ \\
  \hline
  5 & CGL & $(2, 6, 11)^T$ & OPS & $(14, 15, 18)^T$ \\
  \hline
  6 & RPX & $(17, 15, 23)^T$ & USP & $(20, 18, 15)^T$ \\
  \hline
  7 & OUA & $(14, 20, 0)^T$ & ICI & $(8, 2, 8)^T$ \\
  \hline
  8 & IOE & $(8, 14, 4)^T$ & OUS & $(14, 20, 18)^T$ \\
  \hline
  9 & KON & $(10, 14, 13)^T$ & SHA & $(18, 7, 0)^T$ \\
  \hline
  10 & MNA & $(12, 13, 0)^T$ & DYC & $(3, 24, 2)^T$ \\
  \hline
  11 & NIN & $(13, 8, 13)^T$ & HAR & $(7, 0, 17)^T$ \\
  \hline
  12 & LBT & $(11, 1, 19)^T$ & ACT & $(0, 2, 19)^T$ \\
  \hline
  13 & MWT & $(12, 22, 19)^T$ & ERI & $(4, 17, 8)^T$ \\
  \hline
  14 & JFA & $(9, 5, 0)^T$ & SFR & $(18, 5, 17)^T$ \\
  \hline
  15 & CQM & $(2, 16, 12)^T$ & EQU & $(4, 16, 20)^T$ \\
  \hline
  16 & NZA & $(13, 25, 0)^T$ & ENT & $(4, 13, 19)^T$ \\
  \hline
  17 & UIZ & $(20, 8, 25)^T$ & ING & $(8, 13, 6)^T$ \\
  \hline
  18 & IRR & $(8, 17, 17)^T$ & PHA & $(15, 7, 0)^T$ \\
  \hline
  19 & YXQ & $(24, 23, 16)^T$ & RMA & $(17, 12, 0)^T$ \\
  \hline
  20 & OEC & $(14, 4, 2)^T$ & CEU & $(2, 4, 20)^T$ \\
  \hline
  21 & ANI & $(0, 13, 8)^T$ & TIC & $(19, 8, 2)^T$ \\
  \hline
  22 & NLY & $(13, 11, 24)^T$ & ALF & $(0, 11, 5)^T$ \\
  \hline
  23 & ACC & $(0, 2, 2)^T$ & ACI & $(0, 2, 8)^T$ \\
  \hline
  24 & DKP & $(3, 10, 15)^T$ & LIT & $(11, 8, 19)^T$ \\
  \hline
  25 & KIC & $(10, 8, 2)^T$ & YWE & $(24, 22, 4)^T$ \\
  \hline
  26 & RKF & $(17, 10, 5)^T$ & HAV & $(7, 0, 21)^T$ \\
  \hline
  27 & GAZ & $(6, 0, 25)^T$ & ELE & $(4, 11, 4)^T$ \\
  \hline
  28 & XJK & $(23, 9, 10)^T$ & ARN & $(0, 17, 13)^T$ \\
  \hline
  29 & BJO & $(1, 9, 14)^T$ & EDT & $(4, 3, 19)^T$ \\
  \hline
  30 & PWJ & $(15, 22, 9)^T$ & HAT & $(7, 0, 19)^T$ \\
  \hline
  31 & CAI & $(2, 0, 8)^T$ & AMY & $(0, 12, 24)^T$ \\
  \hline
  32 & SEJ & $(18, 4, 9)^T$ & STE & $(18, 19, 4)^T$ \\
  \hline
  33 & CVE & $(2, 21, 4)^T$ & RIO & $(17, 8, 14)^T$ \\
  \hline
  34 & ECK & $(4, 2, 10)^T$ & USC & $(20, 18, 2)^T$ \\
  \hline
  35 & NIN & $(13, 8, 13)^T$ & HAR & $(7, 0, 17)^T$ \\
  \hline
  36 & LBT & $(11, 1, 19)^T$ & ACT & $(0, 2, 19)^T$ \\
  \hline
  37 & LPI & $(11, 15, 8)^T$ & ERH & $(4, 17, 7)^T$ \\
  \hline
  38 & HPR & $(7, 15, 17)^T$ & ASB & $(0, 18, 1)^T$ \\
  \hline
  39 & RFJ & $(17, 5, 9)^T$ & EEN & $(4, 4, 13)^T$ \\
  \hline
  40 & ATI & $(0, 19, 8)^T$ & VIS & $(21, 8, 18)^T$ \\
  \hline
  41 & YQX & $(24, 16, 23)^T$ & ITI & $(8, 19, 8)^T$ \\
  \hline
  42 & GJU & $(6, 9, 20)^T$ & NGE & $(13, 6, 4)^T$ \\
  \hline
  43 & LBV & $(11, 1, 21)^T$ & IEN & $(8, 4, 13)^T$ \\
  \hline
  44 & WHM & $(22, 7, 12)^T$ & TEI & $(19, 4, 8)^T$ \\
  \hline
  45 & BDP & $(1, 3, 15)^T$ & GEN & $(6, 4, 13)^T$ \\
  \hline
  46 & SME & $(18, 12, 4)^T$ & SOK & $(18, 14, 10)^T$ \\
  \hline
  47 & EKG & $(4, 10, 6)^T$ & YOS & $(24, 14, 18)^T$ \\
  \hline
  48 & DLU & $(3, 11, 20)^T$ & ONL & $(14, 13, 11)^T$ \\
  \hline
  49 & VBA & $(21, 1, 0)^T$ & YDR & $(24, 3, 17)^T$ \\
  \hline
  50 & MUI & $(12, 20, 8)^T$ & UGO & $(20, 6, 14)^T$ \\
  \hline
  51 & WDO & $(22, 3, 14)^T$ & RGA & $(17, 6, 0)^T$ \\
  \hline
  52 & TYJ & $(19, 24, 9)^T$ & NIZ & $(13, 8, 25)^T$ \\
  \hline
  53 & KYZ & $(10, 24, 25)^T$ & ATI & $(0, 19, 8)^T$ \\
  \hline
  54 & KIT & $(10, 8, 19)^T$ & ONS & $(14, 13, 18)^T$ \\
  \hline
  55 & JKZ & $(9, 10, 25)^T$ & HED & $(7, 4, 3)^T$ \\
  \hline
  56 & JBN & $(9, 1, 13)^T$ & ISP & $(8, 18, 15)^T$ \\
  \hline
  57 & OBY & $(14, 1, 24)^T$ & LAY & $(11, 0, 24)^T$ \\
  \hline
  58 & SQV & $(18, 16, 21)^T$ & SFA & $(18, 5, 0)^T$ \\
  \hline
  59 & SSX & $(18, 18, 23)^T$ & SHI & $(18, 7, 8)^T$ \\
  \hline
  60 & MWP & $(12, 22, 15)^T$ & ONU & $(14, 13, 20)^T$ \\
  \hline
  61 & GFQ & $(6, 5, 16)^T$ & NCO & $(13, 2, 14)^T$ \\
  \hline
  62 & AWM & $(0, 22, 12)^T$ & MMO & $(12, 12, 14)^T$ \\
  \hline
  63 & QBN & $(16, 1, 13)^T$ & NTO & $(13, 19, 14)^T$ \\
  \hline
  64 & BDP & $(1, 3, 15)^T$ & GEN & $(6, 4, 13)^T$ \\
  \hline
  65 & SME & $(18, 12, 4)^T$ & SOK & $(18, 14, 10)^T$ \\
  \hline
  66 & IMY & $(8, 12, 24)^T$ & YOW & $(24, 14, 22)^T$ \\
  \hline
  67 & XJM & $(23, 9, 12)^T$ & ITH & $(8, 19, 7)^T$ \\
  \hline
  68 & WOL & $(22, 14, 11)^T$ & ADI & $(0, 3, 8)^T$ \\
  \hline
  69 & WGB & $(22, 6, 1)^T$ & STI & $(18, 19, 8)^T$ \\
  \hline
  70 & LOT & $(11, 14, 19)^T$ & NCT & $(13, 2, 19)^T$ \\
  \hline
  71 & MMX & $(12, 12, 23)^T$ & IVE & $(8, 21, 4)^T$ \\
  \hline
  72 & TNJ & $(19, 13, 9)^T$ & SIN & $(18, 8, 13)^T$ \\
  \hline
  73 & QYF & $(16, 24, 5)^T$ & GLE & $(6, 11, 4)^T$ \\
  \hline
  74 & NJV & $(13, 9, 21)^T$ & WIN & $(22, 8, 13)^T$ \\
  \hline
  75 & RLJ & $(17, 11, 9)^T$ & GED & $(6, 4, 3)^T$ \\
  \hline
  76 & HZO & $(7, 25, 14)^T$ & APP & $(0, 15, 15)^T$ \\
  \hline
  77 & LLR & $(11, 11, 17)^T$ & EAR & $(4, 0, 17)^T$ \\
  \hline
  78 & COJ & $(2, 14, 9)^T$ & ANC & $(0, 13, 2)^T$ \\
  \hline
  79 & HJZ & $(7, 9, 25)^T$ & EAN & $(4, 0, 13)^T$ \\
  \hline
  80 & ADH & $(0, 3, 7)^T$ & DHA & $(3, 7, 0)^T$ \\
  \hline
  81 & ACN & $(0, 2, 13)^T$ & SNO & $(18, 13, 14)^T$ \\
  \hline
  82 & EVT & $(4, 21, 19)^T$ & TBE & $(19, 1, 4)^T$ \\
  \hline
  83 & MSP & $(12, 18, 15)^T$ & ENS & $(4, 13, 18)^T$ \\
  \hline
  84 & RFJ & $(17, 5, 9)^T$ & EEN & $(4, 4, 13)^T$ \\
  \hline
  85 & UIZ & $(20, 8, 25)^T$ & ING & $(8, 13, 6)^T$ \\
  \hline
  86 & MSP & $(12, 18, 15)^T$ & ENS & $(4, 13, 18)^T$ \\
  \hline
  87 & COG & $(2, 14, 6)^T$ & OKY & $(14, 10, 24)^T$ \\
  \hline
  88 & SAR & $(18, 0, 17)^T$ & OBE & $(14, 1, 4)^T$ \\
  \hline
  89 & ZWD & $(25, 22, 3)^T$ & FOR & $(5, 14, 17)^T$ \\
  \hline
  90 & DNA & $(3, 13, 0)^T$ & ETH & $(4, 19, 7)^T$ \\
  \hline
  91 & IUB & $(8, 20, 1)^T$ & ERE & $(4, 17, 4)^T$ \\
  \hline
  92 & MWU & $(12, 22, 20)^T$ & ISS & $(8, 18, 18)^T$ \\
  \hline
  93 & KEX & $(10, 4, 23)^T$ & URE & $(20, 17, 4)^T$ \\
  \hline
  94 & LEP & $(11, 4, 15)^T$ & LYN & $(11, 24, 13)^T$ \\
  \hline
  95 & IQG & $(8, 16, 6)^T$ & OWA & $(14, 22, 0)^T$ \\
  \hline
  96 & ZTI & $(25, 19, 8)^T$ & YTH & $(24, 19, 7)^T$ \\
  \hline
  97 & CEO & $(2, 4, 14)^T$ & ISI & $(8, 18, 8)^T$ \\
  \hline
  98 & WWR & $(22, 22, 17)^T$ & SJU & $(18, 9, 20)^T$ \\
  \hline
  99 & OCR & $(14, 2, 17)^T$ & STA & $(18, 19, 0)^T$ \\
  \hline
  100 & NOB & $(13, 14, 1)^T$ & NOR & $(13, 14, 17)^T$ \\
  \hline
  101 & JCD & $(9, 2, 3)^T$ & DIN & $(3, 8, 13)^T$ \\
  \hline
  102 & IIB & $(8, 8, 1)^T$ & ARY & $(0, 17, 24)^T$ \\
  \hline
  103 & DOH & $(3, 14, 7)^T$ & PAT & $(15, 0, 19)^T$ \\
  \hline
  104 & LBV & $(11, 1, 21)^T$ & IEN & $(8, 4, 13)^T$ \\
  \hline
  105 & LMZ & $(11, 12, 25)^T$ & TIN & $(19, 8, 13)^T$ \\
  \hline
  106 & MOJ & $(12, 14, 9)^T$ & WHI & $(22, 7, 8)^T$ \\
  \hline
  107 & KSN & $(10, 18, 13)^T$ & CHC & $(2, 7, 2)^T$ \\
  \hline
  108 & KKY & $(10, 10, 24)^T$ & ASE & $(0, 18, 4)^T$ \\
  \hline
  109 & YQX & $(24, 16, 23)^T$ & ITI & $(8, 19, 8)^T$ \\
  \hline
  110 & PFL & $(15, 5, 11)^T$ & SCL & $(18, 2, 11)^T$ \\
  \hline
  111 & LLR & $(11, 11, 17)^T$ & EAR & $(4, 0, 17)^T$ \\
  \hline
  112 & CND & $(2, 13, 3)^T$ & THA & $(19, 7, 0)^T$ \\
  \hline
  113 & RMX & $(17, 12, 23)^T$ & TSH & $(19, 18, 7)^T$ \\
  \hline
  114 & SCX & $(18, 2, 23)^T$ & EHA & $(4, 7, 0)^T$ \\
  \hline
  115 & GKG & $(6, 10, 6)^T$ & SSO & $(18, 18, 14)^T$ \\
  \hline
  116 & UYQ & $(20, 24, 16)^T$ & MEC & $(12, 4, 2)^T$ \\
  \hline
  117 & FZQ & $(5, 25, 16)^T$ & ONN & $(14, 13, 13)^T$ \\
  \hline
  118 & FXF & $(5, 23, 5)^T$ & ECT & $(4, 2, 19)^T$ \\
  \hline
  119 & XRH & $(23, 17, 7)^T$ & ION & $(8, 14, 13)^T$ \\
  \hline
  120 & TWP & $(19, 22, 15)^T$ & TOT & $(19, 14, 19)^T$ \\
  \hline
  121 & UJQ & $(20, 9, 16)^T$ & HEO & $(7, 4, 14)^T$ \\
  \hline
  122 & WDO & $(22, 3, 14)^T$ & RGA & $(17, 6, 0)^T$ \\
  \hline
  123 & TYJ & $(19, 24, 9)^T$ & NIZ & $(13, 8, 25)^T$ \\
  \hline
  124 & KYZ & $(10, 24, 25)^T$ & ATI & $(0, 19, 8)^T$ \\
  \hline
  125 & OKL & $(14, 10, 11)^T$ & ONW & $(14, 13, 22)^T$ \\
  \hline
  126 & PWJ & $(15, 22, 9)^T$ & HAT & $(7, 0, 19)^T$ \\
  \hline
  127 & TXP & $(19, 23, 15)^T$ & CON & $(2, 14, 13)^T$ \\
  \hline
  128 & MXG & $(12, 23, 6)^T$ & NEC & $(13, 4, 2)^T$ \\
  \hline
  129 & MTK & $(12, 19, 10)^T$ & TIO & $(19, 8, 14)^T$ \\
  \hline
  130 & CRZ & $(2, 17, 25)^T$ & NDO & $(13, 3, 14)^T$ \\
  \hline
  131 & SAI & $(18, 0, 8)^T$ & ESS & $(4, 18, 18)^T$ \\
  \hline
  132 & NMR & $(13, 12, 17)^T$ & HEH & $(7, 4, 7)^T$ \\
  \hline
  133 & YUZ & $(24, 20, 25)^T$ & AVE & $(0, 21, 4)^T$ \\
  \hline
  134 & EVG & $(4, 21, 6)^T$ & TOE & $(19, 14, 4)^T$ \\
  \hline
  135 & LBV & $(11, 1, 21)^T$ & IEN & $(8, 4, 13)^T$ \\
  \hline
  136 & WHM & $(22, 7, 12)^T$ & TEI & $(19, 4, 8)^T$ \\
  \hline
  137 & GPV & $(6, 15, 21)^T$ & THE & $(19, 7, 4)^T$ \\
  \hline
  138 & EFW & $(4, 5, 22)^T$ & REA & $(17, 4, 0)^T$ \\
  \hline
  139 & VUB & $(21, 20, 1)^T$ & RER & $(17, 4, 17)^T$ \\
  \hline
  140 & OOK & $(14, 14, 10)^T$ & UMO & $(20, 12, 14)^T$ \\
  \hline
  141 & ARS & $(0, 17, 18)^T$ & RSA & $(17, 18, 0)^T$ \\
  \hline
  142 & IVY & $(8, 21, 24)^T$ & BOU & $(1, 14, 20)^T$ \\
  \hline
  143 & HAC & $(7, 0, 2)^T$ & NDT & $(13, 3, 19)^T$ \\
  \hline
  144 & PWJ & $(15, 22, 9)^T$ & HAT & $(7, 0, 19)^T$ \\
  \hline
  145 & SQY & $(18, 16, 24)^T$ & EIE & $(4, 8, 4)^T$ \\
  \hline
  146 & GJH & $(6, 9, 7)^T$ & NTE & $(13, 19, 4)^T$ \\
  \hline
  147 & ASH & $(0, 18, 7)^T$ & IHO & $(8, 7, 14)^T$ \\
  \hline
  148 & GQG & $(6, 16, 6)^T$ & USE & $(20, 18, 4)^T$ \\
  \hline
  149 & OUJ & $(14, 20, 9)^T$ & SLU & $(18, 11, 20)^T$ \\
  \hline
  150 & RCF & $(17, 2, 5)^T$ & NAR & $(13, 0, 17)^T$ \\
  \hline
  151 & BFL & $(1, 5, 11)^T$ & IAN & $(8, 0, 13)^T$ \\
  \hline
  152 & JJR & $(9, 9, 17)^T$ & SWH & $(18, 22, 7)^T$ \\
  \hline
  153 & JTA & $(9, 19, 0)^T$ & OFL & $(14, 5, 11)^T$ \\
  \hline
  154 & NPQ & $(13, 15, 16)^T$ & EDF & $(4, 3, 5)^T$ \\
  \hline
  155 & CBK & $(2, 1, 10)^T$ & ROM & $(17, 14, 12)^T$ \\
  \hline
  156 & GPV & $(6, 15, 21)^T$ & THE & $(19, 7, 4)^T$ \\
  \hline
  157 & SUE & $(18, 20, 4)^T$ & MOO & $(12, 14, 14)^T$ \\
  \hline
  158 & NOB & $(13, 14, 1)^T$ & NOR & $(13, 14, 17)^T$ \\
  \hline
  159 & CND & $(2, 13, 3)^T$ & THA & $(19, 7, 0)^T$ \\
  \hline
  160 & RCN & $(17, 2, 13)^T$ & TIT & $(19, 8, 19)^T$ \\
  \hline
  161 & UAE & $(20, 0, 4)^T$ & ISA & $(8, 18, 0)^T$ \\
  \hline
  162 & YSI & $(24, 18, 8)^T$ & SEC & $(18, 4, 2)^T$ \\
  \hline
  163 & XIX & $(23, 8, 23)^T$ & RET & $(17, 4, 19)^T$ \\
  \hline
  164 & NFV & $(13, 5, 21)^T$ & MIL & $(12, 8, 11)^T$ \\
  \hline
  165 & QMN & $(16, 12, 13)^T$ & ITA & $(8, 19, 0)^T$ \\
  \hline
  166 & DSF & $(3, 18, 5)^T$ & RYB & $(17, 24, 1)^T$ \\
  \hline
  167 & KKY & $(10, 10, 24)^T$ & ASE & $(0, 18, 4)^T$ \\
  \hline
  168 & ZWD & $(25, 22, 3)^T$ & FOR & $(5, 14, 17)^T$ \\
  \hline
  169 & GPV & $(6, 15, 21)^T$ & THE & $(19, 7, 4)^T$ \\
  \hline
  170 & BIF & $(1, 8, 5)^T$ & LUN & $(11, 20, 13)^T$ \\
  \hline
  171 & MKT & $(12, 10, 19)^T$ & ARC & $(0, 17, 2)^T$ \\
  \hline
  172 & ACP & $(0, 2, 15)^T$ & API & $(0, 15, 8)^T$ \\
  \hline
  173 & PGJ & $(15, 6, 9)^T$ & TAL & $(19, 0, 11)^T$ \\
  \hline
  174 & YCJ & $(24, 2, 9)^T$ & IFE & $(8, 5, 4)^T$ \\
  \hline
  175 & XJM & $(23, 9, 12)^T$ & ITH & $(8, 19, 7)^T$ \\
  \hline
  176 & SMH & $(18, 12, 7)^T$ & ERO & $(4, 17, 14)^T$ \\
  \hline
  177 & VMQ & $(21, 12, 16)^T$ & FTH & $(5, 19, 7)^T$ \\
  \hline
  178 & EGK & $(4, 6, 10)^T$ & ESE & $(4, 18, 4)^T$ \\
  \hline
  179 & OYP & $(14, 24, 15)^T$ & ARE & $(0, 17, 4)^T$ \\
  \hline
  180 & INB & $(8, 13, 1)^T$ & TRU & $(19, 17, 20)^T$ \\
  \hline
  181 & DNA & $(3, 13, 0)^T$ & ETH & $(4, 19, 7)^T$ \\
  \hline
  182 & DHU & $(3, 7, 20)^T$ & ENJ & $(4, 13, 9)^T$ \\
  \hline
  183 & VRP & $(21, 17, 15)^T$ & UST & $(20, 18, 19)^T$ \\
  \hline
  184 & EKZ & $(4, 10, 25)^T$ & WHA & $(22, 7, 0)^T$ \\
  \hline
  185 & QVC & $(16, 21, 2)^T$ & TIS & $(19, 8, 18)^T$ \\
  \hline
  186 & KRN & $(10, 17, 13)^T$ & THI & $(19, 7, 8)^T$ \\
  \hline
  187 & IIZ & $(8, 8, 25)^T$ & SPE & $(18, 15, 4)^T$ \\
  \hline
  188 & OLQ & $(14, 11, 16)^T$ & RSO & $(17, 18, 14)^T$ \\
  \hline
  189 & SZT & $(18, 25, 19)^T$ & NDE & $(13, 3, 4)^T$ \\
  \hline
  190 & HUD & $(7, 20, 3)^T$ & PEN & $(15, 4, 13)^T$ \\
  \hline
  191 & JCD & $(9, 2, 3)^T$ & DIN & $(3, 8, 13)^T$ \\
  \hline
  192 & NTB & $(13, 19, 1)^T$ & GON & $(6, 14, 13)^T$ \\
  \hline
  193 & ODM & $(14, 3, 12)^T$ & HOW & $(7, 14, 22)^T$ \\
  \hline
  194 & GPV & $(6, 15, 21)^T$ & THE & $(19, 7, 4)^T$ \\
  \hline
  195 & ZDX & $(25, 3, 23)^T$ & SIT & $(18, 8, 19)^T$ \\
  \hline
  196 & PJJ & $(15, 9, 9)^T$ & UAT & $(20, 0, 19)^T$ \\
  \hline
  197 & XRH & $(23, 17, 7)^T$ & ION & $(8, 14, 13)^T$ \\
  \hline
  198 & HKD & $(7, 10, 3)^T$ & DEV & $(3, 4, 21)^T$ \\
  \hline
  199 & QSF & $(16, 18, 5)^T$ & ELO & $(4, 11, 14)^T$ \\
  \hline
  200 & YXW & $(24, 23, 22)^T$ & PSI & $(15, 18, 8)^T$ \\
  \hline
  201 & KZI & $(10, 25, 8)^T$ & TCO & $(19, 2, 14)^T$ \\
  \hline
  202 & HDK & $(7, 3, 10)^T$ & ULD & $(20, 11, 3)^T$ \\
  \hline
  203 & TCV & $(19, 2, 21)^T$ & TUR & $(19, 20, 17)^T$ \\
  \hline
  204 & QJI & $(16, 9, 8)^T$ & NOU & $(13, 14, 20)^T$ \\
  \hline
  205 & UVF & $(20, 21, 5)^T$ & TTO & $(19, 19, 14)^T$ \\
  \hline
  206 & CVA & $(2, 21, 0)^T$ & BEA & $(1, 4, 0)^T$ \\
  \hline
  207 & SAS & $(18, 0, 18)^T$ & SCO & $(18, 2, 14)^T$ \\
  \hline
  208 & KKV & $(10, 10, 21)^T$ & OPA & $(14, 15, 0)^T$ \\
  \hline
  209 & IVY & $(8, 21, 24)^T$ & BOU & $(1, 14, 20)^T$ \\
  \hline
  210 & JYD & $(9, 24, 3)^T$ & TIL & $(19, 8, 11)^T$ \\
  \hline
  211 & GBS & $(6, 1, 18)^T$ & LEG & $(11, 4, 6)^T$ \\
  \hline
  212 & LXC & $(11, 23, 2)^T$ & ALD & $(0, 11, 3)^T$ \\
  \hline
  213 & YFY & $(24, 5, 24)^T$ & RUG & $(17, 20, 6)^T$ \\
  \hline
  214 & MWD & $(12, 22, 3)^T$ & SBE & $(18, 1, 4)^T$ \\
  \hline
  215 & UIZ & $(20, 8, 25)^T$ & ING & $(8, 13, 6)^T$ \\
  \hline
  216 & VBX & $(21, 1, 23)^T$ & MAN & $(12, 0, 13)^T$ \\
  \hline
  217 & COB & $(2, 14, 1)^T$ & UFA & $(20, 5, 0)^T$ \\
  \hline
  218 & GCH & $(6, 2, 7)^T$ & CTU & $(2, 19, 20)^T$ \\
  \hline
  219 & HAD & $(7, 0, 3)^T$ & RED & $(17, 4, 3)^T$ \\
  \hline
  220 & GWH & $(6, 22, 7)^T$ & ATE & $(0, 19, 4)^T$ \\
  \hline
  221 & LBV & $(11, 1, 21)^T$ & IEN & $(8, 4, 13)^T$ \\
  \hline
  222 & WHM & $(22, 7, 12)^T$ & TEI & $(19, 4, 8)^T$ \\
  \hline
  223 & UED & $(20, 4, 3)^T$ & ORS & $(14, 17, 18)^T$ \\
  \hline
  224 & ACM & $(0, 2, 12)^T$ & OME & $(14, 12, 4)^T$ \\
  \hline
  225 & KRN & $(10, 17, 13)^T$ & THI & $(19, 7, 8)^T$ \\
  \hline
  226 & ONE & $(14, 13, 4)^T$ & NGM & $(13, 6, 12)^T$ \\
  \hline
  227 & VBZ & $(21, 1, 25)^T$ & UCH & $(20, 2, 7)^T$ \\
  \hline
  228 & SHY & $(18, 7, 24)^T$ & BIG & $(1, 8, 6)^T$ \\
  \hline
  229 & FFH & $(5, 5, 7)^T$ & GER & $(6, 4, 17)^T$ \\
  \hline
  230 & MYG & $(12, 24, 6)^T$ & WEW & $(22, 4, 22)^T$ \\
  \hline
  231 & HTK & $(7, 19, 10)^T$ & ILL & $(8, 11, 11)^T$ \\
  \hline
  232 & TXP & $(19, 23, 15)^T$ & CON & $(2, 14, 13)^T$ \\
  \hline
  233 & LMZ & $(11, 12, 25)^T$ & TIN & $(19, 8, 13)^T$ \\
  \hline
  234 & ZFT & $(25, 5, 19)^T$ & UET & $(20, 4, 19)^T$ \\
  \hline
  235 & UGS & $(20, 6, 18)^T$ & OGA & $(14, 6, 0)^T$ \\
  \hline
  236 & GPV & $(6, 15, 21)^T$ & THE & $(19, 7, 4)^T$ \\
  \hline
  237 & BOT & $(1, 14, 19)^T$ & RIN & $(17, 8, 13)^T$ \\
  \hline
  238 & ZWD & $(25, 22, 3)^T$ & FOR & $(5, 14, 17)^T$ \\
  \hline
  239 & BRL & $(1, 17, 11)^T$ & MAT & $(12, 0, 19)^T$ \\
  \hline
  240 & XRH & $(23, 17, 7)^T$ & ION & $(8, 14, 13)^T$ \\
  \hline
\end{longtable}


\subsection*{Artefak}

Skrip \textit{solver} yang mengimplementasikan pencarian invers matriks secara modular dengan \textit{Known Plaintext Attack} (KPA) pada bahasa pemrograman Python (\texttt{src/25/solver.py}). Skrip ini secara otomatis memecah teks menjadi blok $m=3$ dan melakukan proses matriks balikan (\textit{decryption}). Log penyelesaian (metadata kunci dan \textit{plaintext}) disimpan di \texttt{src/25/log}.

\subsection*{\textit{Plaintext} Hasil Dekripsi}

\ciphertextfile{../../src/25/log/plain/2026-09-17_201130.txt}

\subsection*{\textit{Plaintext} dengan Format Asli}

\begin{tcolorbox}[breakable, colback=white, colframe=black!80,
  boxrule=0.6pt, arc=0mm, left=3mm, right=3mm, top=2mm, bottom=2mm]
\ttfamily\small\raggedright
Exclusive Scoop!\par
SUSPICIOUS SHADY CHARACTER\par
is frequenting pharmaceutical facility\par
\medskip
We have learned that a mysterious character has been visiting Eientei, Gensokyo\textquotesingle{}s only drug organization. She displays fashion uncommon to Gensokyo and a distinctive single-{}winged appearance, and has not been seen in Gensokyo before. There\textquotesingle{}s surely no way this is just an ordinary patient, in which case it is clear that she has some connection to the organization.\par
\medskip
What connection does she have to Eientei? There are rumors abound that Eientei houses Lunarians who fled from the moon or that it is a secret military base for the Lunar Capital. If either of these are true, then just what is this person? Depending on how the situation develops, it could turn out to be a scoop about illegal drugs being manufactured at Eientei, or something much bigger. We will continue to gather information.\par
\bigskip
Sumber:\par
\hspace*{1.5em}- \url{https://en.touhouwiki.net/wiki/Alternative_Facts_in_Eastern_Utopia/Article_and_Interview/Sagume_Kishin}
\end{tcolorbox}


\subsection*{Kunci}

Matriks Kunci ($K$):
\begin{equation*}
\begin{pmatrix} 18 & 22 & 1 \\ 25 & 12 & 7 \\ 16 & 9 & 11 \end{pmatrix}
\end{equation*}

\subsection*{Referensi}

Konsep \textit{Hill Cipher} dan \textit{known-plaintext attack} mengacu pada bahan kuliah \cite{munir-klasik2}
dan \cite{hill-gfg}. Perhitungan invers matriks modulo 26 mengacu pada materi invers modulo \cite{munir-matematika}.
